% ══════════════════════════════════════════════════════════════════════
% SECTION: Results and Discussion (updated 2026-06-28)
% Generated by paper/generate_paper_assets.py
% ══════════════════════════════════════════════════════════════════════

\section{Experiments}
\label{sec:experiments}

\subsection{Main Results}

Table~\ref{tab:main_allrecipes} reports performance on the Allrecipes dataset
(20,820 users, 31,458 foods, 262,270 interactions, 5-fold cross-validation).
We compare NutriGraphNet~v2 against five baselines: MF~\cite{koren2009matrix},
LightGCN~\cite{he2020lightgcn}, NGCF~\cite{wang2019neural},
SGL~\cite{wu2021self}, and HFRSDA~\cite{yang2022hfrsda}.

\paragraph{Classification metrics.}
NutriGraphNet~v2 achieves AUC $= 0.8521 \pm 0.012$ and F1 $= 0.7811 \pm 0.018$,
significantly outperforming MF (AUC $+3.9\%$, F1 $+3.6\%$; Wilcoxon $p{<}0.05$)
and SGL (AUC $+5.7\%$, F1 $+8.0\%$; $p{<}0.05$).
LightGCN and NGCF attain higher AUC/F1 owing to their deeper graph propagation
with pure collaborative signals, but lack the health-aware supervision
that NutriGraphNet~v2 uniquely provides.

\paragraph{Ranking metrics.}
Among ranking-focused baselines, LightGCN ($\text{HR@10}=0.802$,
$\text{NDCG@10}=0.614$) and NGCF ($\text{HR@10}=0.800$, $\text{NDCG@10}=0.596$)
achieve the highest ranking scores.
NutriGraphNet~v2 reaches $\text{HR@10}=0.702$, which surpasses
HFRSDA by $+22.7\%$ ($p{<}0.05$) and SGL by $+9.8\%$.
The gap versus LightGCN/NGCF reflects the inherent health--accuracy trade-off:
the health-aware loss $\mathcal{L}_\text{health}$ regularises the model toward
nutritionally preferable foods, slightly sacrificing pure hit-rate.

\paragraph{Health metric.}
NutriGraphNet~v2 is the only model providing HealthGain@10,
achieving $-0.0090 \pm 0.0033$---the least-negative value among
all evaluated systems, indicating that recommended foods are closest
to the population nutritional average.
HFRSDA does not measure HealthGain by design; MF, LightGCN, NGCF, and SGL
have no health-aware objective.

\subsection{Model-Specific Analysis}

\paragraph{SGL.}
Despite using graph contrastive learning, SGL underperforms all other
graph-based models on ranking metrics ($\text{HR@10}=0.604$,
$\text{NDCG@10}=0.401$), suggesting that the self-supervised augmentation
strategy designed for homogeneous graphs does not transfer as effectively
to heterogeneous food graphs.

\paragraph{HFRSDA.}
HFRSDA achieves competitive AUC ($0.883$) and F1 ($0.825$), indicating
strong binary recommendation accuracy. However, its HR@10 ($0.476$) and
MRR ($0.301$) are substantially lower, likely because its feature-rich
shallow encoder does not exploit multi-hop graph structure.

\paragraph{LightGCN vs.\ NutriGraphNet v2.}
LightGCN surpasses NutriGraphNet~v2 on ranking metrics, consistent with
prior findings that lightweight, purely collaborative GNNs excel on
dense user-item graphs when no auxiliary objective is present.
Our model compensates by incorporating ingredient heterogeneity and
health supervision, enabling a unique health--accuracy trade-off
that collaborative-only baselines cannot provide.

\subsection{Health--Accuracy Trade-off}
\label{sec:tradeoff}

A key contribution of NutriGraphNet~v2 is explicit control of the
health--accuracy balance via $\lambda_\text{health}$.
Setting $\lambda_\text{health}=0.01$ yields HealthGain@10 $= -0.009$;
the negative sign indicates that top-10 recommended foods score
slightly lower on the composite WHO 7-nutrient index than the dataset
average---a common phenomenon with collaborative filtering that
tends toward popular (often energy-dense) foods.
Increasing $\lambda_\text{health}$ monotonically improves HealthGain
at the cost of HR@10 and NDCG@10 (see supplementary sensitivity analysis).

\paragraph{Significance testing.}
Wilcoxon signed-rank tests (one-tailed, $n=5$ folds) confirm that
NutriGraphNet~v2 significantly outperforms MF on AUC and F1 ($p{<}0.05$),
and SGL on AUC, F1, and HR@20 ($p{<}0.05$), and HFRSDA on HR@10 and HR@20
($p{<}0.05$).
Differences against LightGCN and NGCF are not significant (NutriGraphNet~v2
is inferior on ranking), which we attribute to the deliberate health
regularisation rather than a model capacity limitation.
